% !TEX root = main.tex
\chapter{Data}\label{ch:data}

\section{The corpus}

The corpus holds 1,239 Macedonian poems and song lyrics by 38 authors,
130,059 whitespace-separated tokens in all. Every poem has four fields: a
numeric identifier, the author, the title and the text, in Cyrillic script
with one verse per line and blank lines between stanzas. A typical poem is
short: the median poem has 17 lines and 81 tokens.

The authors are very unevenly represented (Table~\ref{tab:authors}). Three
of them, Блаже Конески, Пијан Славеј and Ефтим Клетников, account for 476
poems, 38\% of the corpus. The next author has 70. Half of the 38 authors
have 17 poems or fewer, and eight have only one or two. Only 24 authors have
at least five poems; these are the authors for whom a stylometric profile
is built (Chapter~\ref{ch:stylometry}), and the remaining 14 appear in the
data only as poems, never as targets of style transfer.

\begin{table}[ht]
\centering
\begin{tabular}{lrrr}
\toprule
Author & Corpus & Train split & Teacher examples \\
\midrule
Блаже Конески       & 182 & 146 & 109 \\
Пијан Славеј        & 160 & 128 & 108 \\
Ефтим Клетников     & 134 & 108 & 110 \\
Славко Јаневски     &  70 &  56 & 120 \\
Влада Урошевиќ      &  61 &  49 & 114 \\
Катица Ќулавкова    &  60 &  48 & 114 \\
Братислав Ташковски &  56 &  44 & 114 \\
Веле Смилевски      &  49 &  39 & 118 \\
Ристо Лазаров       &  48 &  38 & 120 \\
Адем Гајтани        &  45 &  37 & 119 \\
\midrule
Other 28 authors    & 374 & 294 & --- \\
\bottomrule
\end{tabular}
\caption{Poems per author for the ten target authors: in the whole corpus,
in the training split, and synthetic training examples written by the
teacher model with that author as the target (Chapter~\ref{ch:method}).}
\label{tab:authors}
\end{table}

\section{Cleaning}

\paragraph{Poem identifiers.} An author and a title do not identify a poem:
about fifteen (author, title) pairs are shared by several different poems,
Блаже Конески alone having four poems titled \emph{Песна}. Every poem is
therefore given a stable numeric identifier, and every later stage, from
the splits to the generated outputs, refers to poems by that identifier.

\paragraph{Script artefacts.} 52 poems contain a Latin or Greek letter. In
about 22 of them the letter is an artefact of digitisation: a character
that looks like a Cyrillic one but has a different code point, such as a
Greek \emph{о} standing for the vocative \emph{о}, or words that mix two
scripts, such as \emph{штο} or \emph{ΚАМЕН}. The rest are intentional:
Roman numerals and about thirty poems with whole words in Latin script.
The two cases are kept apart. Artefacts are flagged during tagging and
corrected by hand (below); intentional foreign words are left as they are.

\section{Tagging and corrections}

The corpus is tagged with CLASSLA-Stanza \pending{cite: Ljubešić and Dobrovoljc
2019; Terčon and Ljubešić 2023}, using its Macedonian models, which give for
every token a lemma, a universal part-of-speech tag, a MULTEXT-East
morphosyntactic tag \pending{cite: Erjavec 2012} and morphological features.
A hand check of 70 randomly chosen tokens against their lines put the
accuracy of the part-of-speech tag at about 96\%. The stylometric work uses only content
words (nouns, proper nouns, verbs, adjectives and adverbs): 70,531 tokens.

Three corrections are applied on top of the tagger's output:

\begin{enumerate}
  \item \emph{Hand corrections.} A file of corrections, each keyed on a
        word and the tag judged wrong, is applied first, so that correctly
        tagged occurrences of the same word are never rewritten. A
        correction can change the tag, the lemma, or drop the token (for
        noise such as timestamps). There are currently 90.
  \item \emph{Relativisers.} The tagger labels the relativisers
        \emph{што}, \emph{кога}, \emph{како}, \emph{колку} and \emph{каде}
        as adverbs. Left in, they were the five most frequent ``content''
        words in the corpus, 1,478 tokens or 19\% of all adverbs, and
        distorted every part-of-speech ratio built on adverbs. They carry a
        bare adverb tag without a degree, which real adverbs always have,
        and are filtered on that.
  \item \emph{Verbal adverbs.} Forms in \emph{-јќи} (\emph{барајќи}) are
        left unlemmatised by the tagger. They are re-lemmatised to their
        base verb; where the ending is ambiguous (\emph{велејќи} from
        \emph{вели}, but \emph{знаејќи} from \emph{знае}) the more frequent
        verb lemma in the corpus is chosen.
\end{enumerate}

Every tagging run also writes an audit of suspicious tokens, 1,813 at present,
grouped by reason: the two tag layers disagree on the word
class (675 tokens), a lowercase word is tagged as a proper noun (430), the
tagger gives no tag (357, mostly dialectal clippings such as \emph{ко} and
\emph{сал}), a form is tagged differently from its majority reading
elsewhere in the corpus (311), and script artefacts (40). The audit is the
review queue for new hand corrections; nothing is filtered on it directly.
Rules that flagged tokens for being rare were tried and dropped: in poetry
rare words and lines without a verb are normal, and such rules were right
about 5\% of the time.

\section{Splits}

The poems are split into training, validation and test sets in the ratio
80/10/10 within each author, with a fixed random seed: 987, 126 and 126
poems. Splitting by author guarantees that every author with at least three
poems has at least one poem in each set; the eight authors with one or two
poems go entirely to the training set. The split is by poem, so two poems
that share an author and a title can fall into different sets.

The split is computed once and read by every later stage. Everything that
is learned from data (style profiles, classifiers, the teacher's examples,
both students and the reward) uses training poems only. The validation set
is used for model selection and for every comparison made during
development; the test set is used once, for the final numbers in
Chapter~\ref{ch:results}.

\section{How little data there is}

For a task as specific as writing like one poet, this is very little data.
Even the best represented author has 146 poems in the training split, and
seven of the ten target authors have between 37 and 56
(Table~\ref{tab:authors}). By comparison, style-transfer benchmarks such as
\pending{cite: GYAFC (Rao and Tetreault 2018)} have tens of thousands of
parallel sentence pairs for a single, much coarser style distinction.

The synthetic data does not remove this limit. The teacher writes about
115 examples for each target author, a balanced set, but each example is
built from the same training poems: the description of the author's style,
the vocabulary and the example lines in the teacher's prompt all come from
them. An author with 37 poems in the training split is therefore
represented by 37 poems' worth of style, however many examples are written
from it. Chapter~\ref{ch:discussion} returns to this: the students
identifiably imitate only the three authors with more than a hundred
training poems.
